\documentclass[12pt,a4paper]{article}
\usepackage[a4paper,margin=18mm]{geometry}
\usepackage[T2A]{fontenc}
\usepackage[utf8]{inputenc}
\usepackage[russian]{babel}
\usepackage{amsmath,amssymb,mathtools}
\usepackage{array,booktabs,longtable}
\usepackage{xcolor}
\usepackage{hyperref}
\usepackage{tikz}

\hypersetup{
  colorlinks=true,
  linkcolor=blue!55!black,
  urlcolor=blue!55!black
}
\setlength{\parindent}{0pt}
\setlength{\parskip}{5pt}
\setlength{\emergencystretch}{2em}
\renewcommand{\arraystretch}{1.18}

\begin{document}

\begin{titlepage}
\centering
\vspace*{0.18\textheight}
{\LARGE\bfseries Ревью домашней работы\par}
\vspace{18mm}
{\large Автор: Лёвин Артём Александрович\par}
\vspace{6mm}
{\large 16 сентября 2026 года\par}
\vfill
\begin{tikzpicture}
  \draw[line width=0.7pt] (0,0) .. controls (2.2,0.55) and (5.3,-0.55) .. (7.8,0);
\end{tikzpicture}
\end{titlepage}

\newpage
\section*{Сводная таблица}
\small
\begin{longtable}{|p{0.8cm}|p{2.8cm}|p{4.6cm}|p{2.6cm}|p{2.6cm}|}
\hline
\textbf{№} & \textbf{Тема} & \textbf{Суть ошибки} & \textbf{Ваш ответ} & \textbf{Правильный ответ} \\
\hline
\hyperref[task:4]{4} &
Геометрический смысл выражений &
Для $m+2n$ смысл не указан: запись завершена знаком вопроса. &
$mn$ --- площадь; $2m+2n$ --- периметр; $m+2n$ --- ? &
$m+2n$ --- сумма длин трёх сторон: одной стороны длины $m$ и двух сторон ширины $n$. \\
\hline
\hyperref[task:9]{9} &
Минус перед скобками &
При раскрытии $-(x+7)$ знак числа $7$ не изменён: записано $-x+7$ вместо $-x-7$. &
$18$ &
$4$ \\
\hline
\end{longtable}
\normalsize

\newpage
\thispagestyle{empty}
\vspace*{\fill}
\begin{center}
{\LARGE\bfseries Разбор ошибок}
\end{center}
\vspace*{\fill}

\newpage
\phantomsection\label{task:4}
\section*{Задание 4}

\textbf{Текст задания.}

У прямоугольника длина равна $m$, ширина равна $n$. Объясните словами смысл каждого выражения:
\[
mn,\qquad 2m+2n,\qquad m+2n.
\]

\textbf{Ваше решение.}

В записи указано:
\[
mn \text{ --- площадь},
\]
\[
2m+2n \text{ --- периметр},
\]
а для выражения
\[
m+2n
\]
оставлен знак вопроса.

\textbf{Ваш ответ.}

Для первых двух выражений смысл указан, для $m+2n$ ответ не завершён.

\textcolor{red}{Ошибка:} не выполнена обязательная часть задания --- не объяснён геометрический смысл выражения $m+2n$.

\textbf{Где именно возникает ошибка.}

Последний полностью выполненный шаг относится к выражению
\[
2m+2n:
\]
оно верно связано с периметром прямоугольника.

Следующая часть задания требует объяснить выражение
\[
m+2n.
\]
Вместо объяснения после него поставлен знак вопроса. Значит, проблема здесь не в неверном вычислении, а в том, что решение не доведено до конца.

\textbf{Почему это неверно.}

В условии требуется объяснить \emph{каждое} из трёх выражений. Поэтому правильные объяснения первых двух выражений не заменяют третье.

У прямоугольника есть две стороны длины $m$ и две стороны длины $n$. Выражение
\[
m+2n
\]
содержит одну длину $m$ и две ширины $n$. Следовательно, оно показывает суммарную длину трёх сторон прямоугольника: одной стороны длины $m$ и двух сторон длины $n$.

То же самое можно увидеть через периметр. Полный периметр равен
\[
2m+2n.
\]
Если из него убрать одну сторону длины $m$, останется
\[
2m+2n-m=m+2n.
\]
Поэтому $m+2n$ можно также описать как периметр прямоугольника без одной стороны длины $m$.

\textcolor{blue!60!black}{Ключевая идея:} чтобы объяснить геометрический смысл буквенного выражения, нужно сопоставить каждый его член с конкретной длиной, площадью или суммой длин на фигуре. Здесь коэффициент $2$ перед $n$ означает, что берутся две стороны ширины $n$.

\textbf{Исправление от последнего правильного шага.}

Первые два объяснения можно сохранить без изменений:
\[
mn \text{ --- площадь прямоугольника},
\]
\[
2m+2n \text{ --- периметр прямоугольника}.
\]

Нужно дописать только третий пункт:
\[
m+2n
\]
--- это сумма длин трёх сторон прямоугольника: одной стороны длины $m$ и двух сторон длины $n$.

\textbf{Полное правильное решение.}

Площадь прямоугольника равна произведению его длины на ширину, поэтому
\[
mn
\]
--- площадь прямоугольника.

Периметр равен сумме длин всех четырёх сторон. Две стороны имеют длину $m$, ещё две --- длину $n$, поэтому
\[
2m+2n
\]
--- периметр прямоугольника.

В выражении
\[
m+2n
\]
слагаемое $m$ задаёт длину одной стороны длины $m$, а $2n$ --- суммарную длину двух сторон ширины $n$. Следовательно,
\[
m+2n
\]
--- суммарная длина трёх сторон прямоугольника: одной стороны длины $m$ и двух сторон длины $n$.

\textbf{Проверка.}

Возьмём для проверки прямоугольник с длиной $m=8$ и шириной $n=3$. Тогда три стороны, о которых идёт речь, имеют длины
\[
8,\quad 3,\quad 3.
\]
Их сумма равна
\[
8+3+3=14.
\]
По выражению получаем
\[
m+2n=8+2\cdot 3=14.
\]
Значения совпали, значит геометрический смысл сформулирован верно.

\textcolor{teal!60!black}{Итог:} первые два объяснения верны; нужно было дописать, что $m+2n$ --- сумма длин трёх сторон прямоугольника: одной стороны длины $m$ и двух сторон длины $n$.

\textbf{Задача на закрепление.}

У прямоугольника длина равна $a$, ширина равна $b$. Объясните словами геометрический смысл выражений
\[
ab,\qquad 2a+2b,\qquad 2a+b.
\]

\newpage
\phantomsection\label{task:9}
\section*{Задание 9}

\textbf{Текст задания.}

Упростите выражение и затем при $x=4$ найдите его значение:
\[
3(2x-5)-(x+7)+2(x-1).
\]

\textbf{Ваше решение.}

В записи выполнено раскрытие скобок:
\[
6x-15-x+7+2x-2,
\]
после чего получено
\[
7x-10,
\]
а при $x=4$:
\[
28-10=18.
\]

\textbf{Ваш ответ.}

\[
18.
\]

\textcolor{red}{Ошибка:} неверно раскрыта скобка, перед которой стоит знак минус.

\textbf{Где именно возникает ошибка.}

До второй скобки ход решения верный:
\[
3(2x-5)=6x-15.
\]
Также верно раскрыта последняя скобка:
\[
2(x-1)=2x-2.
\]

Первый достоверно неверный шаг возникает при переходе
\[
-(x+7)\longrightarrow -x+7.
\]

Правильно должно быть
\[
-(x+7)=-x-7.
\]

\textbf{Почему этот шаг неверен.}

Минус перед скобками означает умножение \emph{всего} выражения в скобках на $-1$:
\[
-(x+7)=(-1)(x+7).
\]
По распределительному закону число $-1$ нужно умножить на каждое слагаемое:
\[
(-1)\cdot x=-x,
\]
\[
(-1)\cdot 7=-7.
\]
Поэтому оба знака внутри скобок меняются:
\[
-(x+7)=-x-7.
\]

В Вашей записи знак у $x$ изменён правильно, а знак у $7$ оставлен положительным. Из-за этого постоянная часть выражения стала больше на $14$: вместо $-7$ появилось $+7$. Эта ошибка затем перешла и в упрощённое выражение, и в числовой ответ.

\textcolor{blue!60!black}{Ключевая идея:} если перед скобками стоит минус без другого числа, можно считать, что перед скобками стоит множитель $-1$. Он умножается на каждый член скобки, поэтому знак каждого слагаемого меняется на противоположный.

\textbf{Исправление от последнего правильного шага.}

Сохраним верное раскрытие первой скобки и исправим только вторую:
\[
3(2x-5)-(x+7)+2(x-1)
=
6x-15-x-7+2x-2.
\]

Теперь отдельно соберём слагаемые с $x$:
\[
6x-x+2x=7x.
\]

Отдельно сложим числа:
\[
-15-7-2=-24.
\]

Поэтому после упрощения получаем
\[
7x-24.
\]

Теперь подставим $x=4$:
\[
7\cdot 4-24=28-24=4.
\]

\textbf{Полное правильное решение.}

\[
\begin{aligned}
3(2x-5)-(x+7)+2(x-1)
&=6x-15-x-7+2x-2\\
&=(6x-x+2x)+(-15-7-2)\\
&=7x-24.
\end{aligned}
\]

При $x=4$:
\[
7\cdot 4-24=28-24=4.
\]

\textbf{Проверка.}

Проверим значение напрямую в исходном выражении:
\[
3(2\cdot 4-5)-(4+7)+2(4-1).
\]
Сначала вычислим скобки:
\[
2\cdot 4-5=3,\qquad 4+7=11,\qquad 4-1=3.
\]
Тогда
\[
3\cdot 3-11+2\cdot 3=9-11+6=4.
\]
Прямая подстановка дала то же значение, что и упрощённое выражение.

\textcolor{teal!60!black}{Итог:} при раскрытии $-(x+7)$ нужно менять знак у обоих слагаемых. После исправления выражение упрощается до $7x-24$, а при $x=4$ его значение равно $4$.

\textbf{Задача на закрепление.}

Упростите выражение и затем при $x=3$ найдите его значение:
\[
4(2x-3)-(x+5)+3(x-2).
\]

\vfill
\begin{center}
\small Лёвин Артём Александрович эксклюзивно для Кирилла
\end{center}

\end{document}
